\documentclass{beamer}
\usetheme{Madrid}
\usecolortheme{default}

\usepackage[utf8]{inputenc}
\usepackage[T1]{fontenc}
\usepackage{booktabs}
\usepackage{array}
\usepackage{listings}
\usepackage{xcolor}

\title{Geospatial Dataloader Market Survey\\and Comparison with Our V2}
\author{Project Internal Review}
\date{\today}

\lstdefinestyle{code}{
  basicstyle=\ttfamily\scriptsize,
  frame=single,
  breaklines=true,
  showstringspaces=false
}

\begin{document}

\begin{frame}
  \titlepage
\end{frame}

\begin{frame}{Survey Scope}
\textbf{Question:} How do mainstream geospatial libraries build dataloaders, and where does our \texttt{dataloader\_v2} stand?
\vspace{0.5em}
\begin{itemize}
  \item Focus: map/EO data loading, alignment, sampling, label handling, synthetic time filling.
  \item Sources: official docs and GitHub pages (TorchGeo, Raster Vision, eo-learn, stackstac/xarray, Open Data Cube, FarmVibes.AI).
  \item Comparison target: current capabilities in our \texttt{dataloader\_v2}.
\end{itemize}
\end{frame}

\begin{frame}{What the Market Commonly Does}
\begin{itemize}
  \item \textbf{Unify to a tensor/datacube contract}: usually \texttt{[time, band, y, x]} or patch-first tensors.
  \item \textbf{Spatial-temporal query first}: \texttt{bbox/roi + time range + resolution} are the core API knobs.
  \item \textbf{On-the-fly reprojection/resampling}: align mixed CRS/resolution during load, not by rewriting raw files.
  \item \textbf{Patch/chip sampling}: random/grid samplers for training and inference scalability.
  \item \textbf{Lazy/parallel I/O}: Dask/chunked reads for large rasters and cloud assets.
  \item \textbf{Label decoupling}: features and labels can come from different products/sources.
\end{itemize}
\end{frame}

\begin{frame}{Synthetic / Gap Filling in Practice}
\begin{itemize}
  \item Most libraries support \textbf{resample + interpolation} for continuous variables.
  \item Typical methods: \texttt{nearest}, \texttt{ffill/bfill}, linear/cubic interpolation.
  \item Event-like signals are often handled with \textbf{carry-forward} or sparse encoding, not smooth interpolation.
  \item Good practice: expose whether values are observed vs imputed using masks/flags.
  \item Gap-filling is usually optional and task-driven, not always-on by default.
\end{itemize}
\vspace{0.4em}
\textbf{Our V2 matches this direction} with \texttt{synthetic\_time}, method controls, and synthetic masks.
\end{frame}

\begin{frame}{Representative Ecosystem Snapshot}
\small
\begin{tabular}{>{\raggedright\arraybackslash}p{2.2cm} >{\raggedright\arraybackslash}p{2.5cm} >{\raggedright\arraybackslash}p{2.5cm} >{\raggedright\arraybackslash}p{2.8cm}}
\toprule
Library & Loader Pattern & Synthetic/Resample & Typical Positioning \\
\midrule
TorchGeo & GeoDataset + GeoSampler (random/grid/batch) & Not a dedicated synthetic module; relies on transforms/workflows & Deep learning training pipelines on geospatial data \\
\midrule
Raster Vision & Config-driven pipeline + chip generation & Limited built-in interpolation focus & End-to-end CV workflow over large imagery \\
\midrule
eo-learn & EOPatch + task graph & Strong interpolation/resample tasks & EO time-series engineering and preprocessing \\
\midrule
stackstac + xarray & STAC to lazy xarray datacube & Strong resample/interp APIs & Cloud-native analysis and large-scale data fusion \\
\midrule
Open Data Cube & Product-indexed query + load & Resampling/grouping support, pipeline-specific gap handling & National/enterprise datacube operations \\
\bottomrule
\end{tabular}
\end{frame}

\begin{frame}{Our V2: Current Strengths}
\begin{itemize}
  \item Unified simple API: \texttt{data/date\_range/bbox/resolution} with stable outputs \texttt{x[T,C,H,W]}, \texttt{y[T,H,W]}.
  \item Multi-source fusion already wired (FIRMS, ERA5, NOAA, LANDFIRE, MERRA2 where available).
  \item Label source decoupling is implemented (\texttt{mtbs/noaa/firms/None}).
  \item Explicit synthetic controls plus observed-vs-synthetic masks in \texttt{meta}.
  \item Fast source availability screening and user-facing warnings improve reliability.
  \item No raw data mutation required; normalization is done at load time.
\end{itemize}
\end{frame}

\begin{frame}{Our V2: Gaps vs Mature Platforms}
\begin{itemize}
  \item No generalized sampler module yet (random/grid/stratified chips as first-class objects).
  \item Local-file first; limited cloud-native STAC/COG lazy loading compared with stackstac ecosystem.
  \item Synthetic methods are practical but still basic vs richer temporal models.
  \item Metadata/provenance schema can be more standardized (lineage, quality flags, per-source confidence).
  \item Ecosystem integration is thinner (fewer plug-in adapters and benchmark-ready datasets).
\end{itemize}
\end{frame}

\begin{frame}{Decision View: Why V2 Is Still Valuable}
\begin{itemize}
  \item For our current client workflow, V2 is already \textbf{deployable} and task-aligned.
  \item It gives a clean contract for algorithms while preserving source flexibility.
  \item It handles core operational pain points: missing data, temporal mismatch, and label configurability.
  \item It avoids over-engineering before product requirements stabilize.
\end{itemize}
\vspace{0.4em}
\textbf{Conclusion:} V2 is a strong pragmatic baseline; we should now selectively absorb market patterns.
\end{frame}

\begin{frame}{Recommended Next Iteration (V2.x)}
\begin{itemize}
  \item Add a first-class sampler API (\texttt{sampler="grid|random"} + patch size/stride).
  \item Add STAC/COG reader option for lazy cloud/local unified loading.
  \item Add preset synthetic policies by variable type (\texttt{continuous/event/label}) with defaults.
  \item Add run diagnostics report: coverage, missing rate, synthetic ratio per channel/source.
  \item Keep warning UX and source-decoupled labels as core differentiators.
\end{itemize}
\end{frame}

\begin{frame}{References (Web)}
\footnotesize
\begin{itemize}
  \item TorchGeo docs (samplers): \texttt{https://torchgeo.readthedocs.io/en/latest/api/samplers.html}
  \item TorchGeo GitHub: \texttt{https://github.com/microsoft/torchgeo}
  \item Raster Vision: \texttt{https://rastervision.io/}
  \item eo-learn EOTasks: \texttt{https://eo-learn.readthedocs.io/en/latest/eotasks.html}
  \item eo-learn interpolation task docs: \texttt{https://eo-learn.readthedocs.io/en/stable/\_modules/eolearn/features/extra/interpolation.html}
  \item stackstac docs: \texttt{https://stackstac.readthedocs.io/}
  \item xarray time resample/interpolate: \texttt{https://docs.xarray.dev/en/v2025.06.0/user-guide/time-series.html}
  \item Open Data Cube loading: \texttt{https://opendatacube.readthedocs.io/en/1.8.20/api/indexed-data/loading.html}
  \item FarmVibes.AI workflows: \texttt{https://microsoft.github.io/farmvibes-ai/docfiles/markdown/WORKFLOWS.html}
\end{itemize}
\end{frame}

\end{document}
